\chapter[Estado del arte]{Estado del arte: de la recuperación léxica a los modelos de lenguaje}
\label{ch:sota}

El problema planteado (emparejar una descripción libre con el perfil de un experto) se sitúa en la intersección de dos campos bien establecidos: la búsqueda de expertos (\emph{expertise retrieval}) y la \gls{ri} textual.
En lugar de recorrer su evolución de forma cronológica, este capítulo organiza la revisión en torno a las tres decisiones de diseño que estructuran un sistema como el propuesto: \emph{cómo representar al experto} (Secciones~\ref{sec:sota-expertos}-\ref{sec:sota-dispersa}), \emph{cómo representar la consulta del estudiante} y \emph{cómo combinar ambas representaciones} para producir un ranking (Secciones~\ref{sec:sota-fusion} y~\ref{sec:sota-llm}).
Cada técnica se introduce, por tanto, no por su interés histórico, sino por la necesidad concreta del sistema que ayuda a cubrir: interpretabilidad, robustez ante la brecha de vocabulario, multilingüismo o despliegue local.
El capítulo concluye posicionando este trabajo: qué toma de cada familia y qué aporta sobre ellas.

\section{Búsqueda de expertos}
\label{sec:sota-expertos}

La búsqueda de expertos estudia cómo identificar, dentro de una organización, a las personas con conocimiento sobre un tema dado a partir de la evidencia documental que generan.
El área quedó consolidada a partir de las campañas de evaluación empresarial de \gls{trec} y está sistematizada en la revisión de referencia de \citet{balog2012expertise}, que distingue dos familias de modelos: los \emph{centrados en el candidato}, que construyen una representación agregada de cada persona a partir de sus documentos y la comparan con la consulta, y los \emph{centrados en el documento}, que recuperan primero documentos relevantes y trasladan después esa relevancia a las personas asociadas.
Esta monografía sintetiza más de una década de trabajo nacido de dichas campañas y fija la terminología y la taxonomía de modelos que siguen vigentes hoy, lo que la convierte en el anclaje natural para situar nuestra tarea dentro del campo.

Un ejemplo canónico de la aproximación centrada en el candidato es el \emph{Toronto Paper Matching System} \citep{charlin2013tpms}, que en la asignación de revisores construye un perfil agregado de cada experto a partir de sus publicaciones previas y estima su afinidad con cada $texto-consulta$, patrón que este trabajo reutiliza en el dominio de la tutorización.

La recomendación de tutores académicos es una instanciación natural de este marco: la evidencia son publicaciones, proyectos y trabajos supervisados, donde el «experto» buscado es el tutor cuya trayectoria mejor cubre la idea del estudiante.
Este trabajo adopta la aproximación centrada en el candidato (cada investigador se representa mediante perfiles agregados de su actividad) por dos razones: produce un índice compacto sobre el que la búsqueda es inmediata y genera representaciones interpretables (temas, dominios, biografía) que pueden mostrarse al estudiante como justificación de la recomendación.
Esto importa en nuestro problema porque la representación del experto debe ser, a la vez, la unidad sobre la que se busca y la evidencia que justifica la recomendación ante el estudiante (requisito \ref{req:emparejamiento}): la vía centrada en el candidato satisface ambas necesidades de partida.

\section{Recuperación léxica: del modelo vectorial a BM25}
\label{sec:sota-lexica}

La tradición léxica representa consultas y documentos como vectores de términos ponderados y estima la relevancia a partir del solapamiento de esos términos: cuantos más términos importantes comparten una consulta y un documento, más afines se consideran.
Toda la familia descansa, por tanto, en la coincidencia de superficie de las palabras, una premisa tan eficaz cuando consulta y documento comparten vocabulario como frágil cuando no lo hacen.
Sobre el modelo de espacio vectorial de \citet{salton1975vector} y los esquemas de pesado \gls{tfidf}\footnote{La ponderación \gls{tfidf} combina la frecuencia de un término en el documento (\emph{term frequency}) con su rareza en la colección (\emph{inverse document frequency}): pesan más los términos abundantes en un documento pero infrecuentes en el resto, que son los más discriminativos.} se construyó la función de ordenación \gls{bm25} \citep{robertson2009bm25}, que pondera la frecuencia del término saturándola y normalizando por longitud del documento y que sigue siendo (décadas después) una línea base extraordinariamente difícil de batir.

Las fortalezas de esta familia son precisamente las que un sistema institucional agradece: eficiencia mediante índices invertidos, ausencia de entrenamiento y, sobre todo, \emph{interpretabilidad} (es posible mostrar qué términos del perfil han producido la coincidencia).
Su debilidad estructural es el desajuste de vocabulario descrito en la Sección~\ref{sec:problema-vocabulario}: si estudiante e investigador no comparten términos, no hay coincidencia posible, por muy afines que sean los conceptos.
En nuestro contexto, donde la consulta llega en español coloquial y los perfiles están redactados en lenguaje técnico, esta debilidad es determinante y obliga a complementar lo léxico con representaciones semánticas.

\section{Representaciones densas y búsqueda semántica}
\label{sec:sota-densa}

La hipótesis distribucional (el significado de una palabra se deduce de sus contextos) dio lugar primero a los vectores de palabras \citep{mikolov2013word2vec} y, con la arquitectura Transformer \citep{vaswani2017attention} y los modelos preentrenados tipo BERT \citep{devlin2019bert}, a representaciones contextuales de frases y documentos completos.
Los \emph{bi-encoders} como Sentence-BERT \citep{reimers2019sbert} proyectan textos a un \emph{espacio vectorial} en el que la proximidad geométrica aproxima la afinidad semántica, y trabajos como \emph{Dense Passage Retrieval} \citep{karpukhin2020dpr} demostraron que esta recuperación densa puede superar a \gls{bm25} en tareas de respuesta a preguntas.

Para este trabajo, dos propiedades de los \emph{embeddings} densos modernos resultan decisivas.
La primera es su carácter \emph{multilingüe}: modelos abiertos de \emph{embeddings} entrenados sobre corpus en múltiples idiomas \citep{nussbaum2025nomic} proyectan «detección de bulos en redes sociales» y \emph{misinformation detection} a regiones cercanas del espacio, lo que ataca de raíz la doble brecha de idioma y vocabulario (requisito \ref{req:emparejamiento}).
La segunda es su disponibilidad como modelos abiertos ejecutables en local, compatible con el requisito \ref{req:soberania}.
Su contrapartida es la opacidad: la similitud de coseno\footnote{Mide la afinidad entre dos vectores como el coseno del ángulo que forman. Vale 1 cuando apuntan en la misma dirección y es independiente de su magnitud, lo que la convierte en la medida estándar para comparar \emph{embeddings} densos.} no explica por sí misma \emph{por qué} dos textos se parecen, lo que penaliza la transparencia exigible a una herramienta de orientación.

\section{Recuperación dispersa neuronal: lo mejor de ambos mundos}
\label{sec:sota-dispersa}

Una línea intermedia busca conservar la eficiencia e interpretabilidad del índice invertido incorporando aprendizaje neuronal al pesado de términos.
SPLADE \citep{formal2021splade} aprende representaciones dispersas con expansión de términos. Más recientemente, BM42 \citep{vasnetsov2024bm42} propone ponderar los términos de un texto mediante las atenciones de un modelo Transformer, combinándolas con la \gls{idf} al estilo \gls{bm25}.
El resultado es un vector disperso cuyos componentes siguen siendo términos legibles (cada coincidencia es explicable) pero cuyos pesos reflejan la importancia semántica del término en su contexto, no su mera frecuencia.

Este trabajo emplea esta familia para representar los vocabularios temáticos de los investigadores (temas declarados y dominios de especialización extraídos automáticamente), reservando los \emph{embeddings} densos para las descripciones biográficas extensas.

\section{Hibridación y fusión de rankings}
\label{sec:sota-fusion}

Dado que las señales léxicas y semánticas fallan en casos distintos, su combinación es habitualmente superior a cualquiera de ellas por separado.
Combinar una señal de coincidencia de términos o áreas con una señal textual semántica es, de hecho, una práctica consolidada en la asignación de revisores a gran escala, donde la afinidad $experto-documento$ se compone de varias señales heterogéneas \citep{charlin2013tpms}.
La diferencia aquí está en cómo se combinan, pues los esquemas habituales exigen calibrar pesos entre escalas no comparables.
Entre las técnicas de fusión, la \gls{rrf} \citep{cormack2009rrf} destaca por su robustez y simplicidad: combina rankings atendiendo solo a las posiciones, sin necesidad de calibrar puntuaciones heterogéneas,
\begin{equation}
\mathrm{RRF}(d) \;=\; \sum_{r \in R} \frac{1}{k + \mathrm{rank}_r(d)},
\label{eq:rrf}
\end{equation}
donde $R$ es el conjunto de rankings a fusionar, $\mathrm{rank}_r(d)$ la posición del documento $d$ en el ranking $r$, y $k$ una constante de suavizado (típicamente $k=60$).
Su independencia de las escalas de puntuación la hace especialmente adecuada para fusionar una búsqueda dispersa (puntuaciones BM42) con una densa (similitudes de coseno), que viven en rangos no comparables.

\section{Modelos de lenguaje en el bucle de recuperación}
\label{sec:sota-llm}

La irrupción de los \glspl{llm} \citep{brown2020gpt3} ha abierto una tercera vía: utilizarlos no como buscadores, sino como \emph{transformadores de la consulta} o \emph{sintetizadores del corpus} alrededor de un motor de recuperación clásico.
En la dirección $consulta\Rightarrow corpus$, técnicas como Query2doc \citep{wang2023query2doc} e HyDE \citep{gao2023hyde} expanden o reescriben la consulta (incluso generando un documento hipotético) para acercarla al registro del corpus antes de buscar. Trasladado a nuestro problema, un \gls{llm} puede traducir la idea coloquial del estudiante al inglés técnico de los perfiles y enriquecerla con terminología afín, sin que el estudiante deba conocerla.
En la dirección $\Rightarrow consulta$, el patrón de \gls{rag} \citep{lewis2020rag} consolidó la arquitectura $recuperador+generador$. Aquí lo relevante es el uso inverso: emplear el \gls{llm} para \emph{sintetizar} cada trayectoria investigadora en una biografía homogénea y densa en información, que se convierte en la unidad de indexación semántica.
Esta estrategia de resumir historiales ricos en texto para alimentar un recomendador, ponderando la actividad más reciente, se ha explorado en recomendación secuencial \citep{zheng2024harnessing}.
El resumen \emph{recursivo} de textos que exceden la ventana de contexto cuenta, además, con precedentes consolidados en recuperación, como RAPTOR \citep{sarthi2024raptor}, que construye de abajo arriba un árbol de resúmenes para representar documentos extensos a varias granularidades.
En este trabajo se adopta esta idea para construir la biografía de cada investigador de forma incremental, dando más peso a sus años recientes.
La variante propuesta agrupa la producción \emph{cronológicamente}, por periodos, en lugar de por similitud, y conserva el resumen de cada periodo como artefacto trazable (Sección~\ref{sec:met-indexacion}).

En ambos casos el \gls{llm} desempeña un papel \emph{acotado y auxiliar} (enriquecer perfiles y reformular o interpretar consultas), nunca el de generar libremente la respuesta final.
Esta delimitación es deliberada: mantener al modelo fuera del bucle de decisión preserva un sistema medible y auditable, y reduce su dependencia de la generación libre, cuya validación es notoriamente más difícil.

Dos avances recientes hacen viable este enfoque bajo el requisito de despliegue local (\ref{req:soberania}).
El primero es la madurez de los \glspl{llm} abiertos de tamaño medio ejecutables en local, con capacidad multilingüe e instruccional suficiente, como las familias Qwen \citep{qwen2025qwen25}, Gemma, Llama o Mistral.
El segundo son los marcos de \emph{salida estructurada}, que permiten obtener de un \gls{llm} objetos validables (listas de términos, campos tipados) en lugar de texto libre, condición necesaria para integrarlo de forma fiable en un pipeline automatizado y evaluable.

\section{Posicionamiento de este trabajo}
\label{sec:sota-posicionamiento}

Conviene subrayarlo sin ambigüedad: este trabajo \emph{no} propone un nuevo algoritmo de recuperación.
Su contribución es de \emph{integración y medida}, y se concreta en tres planos: (i)~integrar técnicas conocidas, representativas de cada familia del estado del arte, en un sistema institucional real y desplegable en infraestructura propia; (ii)~construir, a partir de datos públicos, perfiles semánticos del profesorado que sirvan de representación del experto; y (iii)~evaluar el conjunto sobre un mismo corpus mediante una verdad de referencia objetiva y reproducible (Capítulo~\ref{ch:evaluacion}).
La Tabla~\ref{tab:posicionamiento} resume la correspondencia entre las estrategias que se comparan y las familias del estado del arte de las que proceden.

\begin{table}[htbp]
    \centering
    \caption[Correspondencia entre estrategias y familias del estado del arte]{Correspondencia entre las estrategias del sistema y las familias del estado del arte.}
    \label{tab:posicionamiento}
    \begin{tabularx}{\textwidth}{l X X}
        \toprule
        \rowcolor{upmblue!8} \colhead{Estrategia del sistema} & \colhead{Descripción} & \colhead{Marco conceptual} \\
        \midrule
        Léxica con \gls{llm} & Extracción de términos de la idea con \gls{llm} y búsqueda dispersa BM42 con fusión RRF & Recuperación dispersa neuronal y \glspl{llm} (§\ref{sec:sota-dispersa}, §\ref{sec:sota-llm}) \\
        Léxica directa & Términos propios de la consulta sobre el índice disperso, sin \gls{llm} (ablación) & Recuperación léxica y dispersa (§\ref{sec:sota-lexica}, §\ref{sec:sota-dispersa}) \\
        Semántica directa & Embedding denso de la idea tal cual y similitud de coseno sobre biografías & Representaciones densas (§\ref{sec:sota-densa}) \\
        Semántica enriquecida & Reescritura y expansión de la idea con \gls{llm} (estilo HyDE/Query2doc) y búsqueda densa & Representaciones densas y \glspl{llm} (§\ref{sec:sota-densa}, §\ref{sec:sota-llm}) \\
        Híbrida & Fusión RRF de la señal dispersa (dominios) y la densa (biografías) & Hibridación y fusión de rankings (§\ref{sec:sota-fusion}) \\
        \bottomrule
    \end{tabularx}
\end{table}

Frente a los trabajos de búsqueda de expertos clásicos, se aporta la incorporación sistemática de \glspl{llm} locales en la construcción del perfil y en la transformación de la consulta.
Frente a las demostraciones recientes basadas en \glspl{llm}, se aporta un marco de evaluación sin etiquetado manual (Capítulo~\ref{ch:evaluacion}) que permite afirmar (y no solo ilustrar) qué estrategia funciona mejor en este dominio.

Conviene, por último, situar el valor de fondo de esta aproximación.
Mientras que los enfoques clásicos de búsqueda de expertos representan a cada investigador mediante listas de palabras clave o vectores derivados exclusivamente de sus publicaciones, este trabajo explora una vía híbrida en la que los modelos generativos contribuyen a construir representaciones semánticas más ricas, interpretables y cercanas al lenguaje humano.
El esfuerzo principal se dirige así a transformar una información académica dispersa y heterogénea en perfiles semánticos estructurados que resultan legibles \emph{tanto} por los algoritmos de recuperación \emph{como} por sistemas de \gls{ia} generativa.
En esa doble legibilidad reside una proyección natural del trabajo más allá del recomendador, que se retoma entre las líneas futuras (Capítulo~\ref{ch:conclusiones}).
\pendiente{Si se identifican sistemas análogos de recomendación de tutores en otras universidades (trabajos publicados o herramientas en producción), añadir aquí una breve revisión comparativa.}
